% Core labelled calculus for K from Negri's geometric-rule method.
\calculussheet{sheet:labelled-k-negri}{Negri Labelled Sequent Calculus for $\mathsf K$}{Language: labelled formulas $x:A$ and relational atoms $xRy$; sequents are finite multisets of both kinds on each side; classical propositional LK rules act at one fixed label; modal rules are displayed below; labels are metalinguistic world variables; equality: label identity only through uniform renaming, no object-language equality.}
\noindent\emph{Controlling source:} \cite{Negri2005}.
\hbrule{rule:labelled-k-id}{Identity}{ }{x:A\Rightarrow x:A}{\none}
\hbrule{rule:labelled-k-imp-right}{Implication-R}{\Gamma,x:A\Rightarrow x:B,\Delta}{\Gamma\Rightarrow x:(A\to B),\Delta}{\none}
\hbrule{rule:labelled-k-imp-left}{Implication-L}{\Gamma\Rightarrow x:A,\Delta\qquad\Gamma,x:B\Rightarrow\Delta}{\Gamma,x:(A\to B)\Rightarrow\Delta}{contexts are combined in the ordinary multiset presentation}
\hbrule{rule:labelled-k-box-right}{Box-R}{xRy,\Gamma\Rightarrow\Delta,y:A}{\Gamma\Rightarrow\Delta,x:\Box A}{$y$ is fresh for the conclusion}
\hbrule{rule:labelled-k-box-left}{Box-L}{xRy,y:A,x:\Box A,\Gamma\Rightarrow\Delta}{xRy,x:\Box A,\Gamma\Rightarrow\Delta}{\none}
